% =============================================================================
% 007 / 068
% Status: DONE
% =============================================================================
% Notes:
%
% Source question:
% 您知道有誰系統地寫了吳越國史嗎？如果沒有，我打算在墻內現擬定一份初稿。因為才疏學淺，期間遇到問題，會繼續在此處向您發問。
% =============================================================================

\chapter[您知道有誰系統地寫了吳越國史嗎？如果沒有，我打算在墻內現擬定一份初稿。因為才疏學淺，期間遇到問題，會繼續在此處向您發問。]{您知道有誰系統地寫了吳越國史嗎？如果沒有，我打算在墻內現擬定一份初稿。因為才疏學淺，期間遇到問題，會繼續在此處向您發問。}
\qaanswer{
吳語公眾號有人在更吧。如果在牆內建議敏感內容還是不要公佈的好。

寫史先得確立歷史觀，構建吳越史採用的歷史觀建議參考UK的歷史構建觀。因為吳越和UK非常類似，同樣都是由少數原住民，歷史上不斷地以武力、半武力方式遷入的移民層層交錯疊合而成的文明體；同樣在在出現穩定的文化、文明特質後（吳越是南宋以後\ruby{󱚾󱤦}{UK}是諾曼征服者進入以後），都沒有經歷過釜底抽薪的大浩劫（戰爭、瘟疫引起的人口滅絕、人口置換），且內部比較好的融合出一套獨具自身特點的不同與周邊地區的語言、文化、風俗及工商產業結構。
}

